High-risk AI systems must now demonstrate conformity with overlapping governance
frameworks, yet the technical evaluation literature and the regulatory requirements have
developed in separate worlds. This article bridged them. We organized the AI-evaluation
literature into twelve method families and mapped each, graded by technical relevance, to
the requirements of the NIST AI RMF, the EU AI Act high-risk obligations, and ISO/IEC
42001---producing, to our knowledge, the first crosswalk of evaluation methods to the
simultaneous requirements of the three frameworks within the technical literature we surveyed. The
resulting matrix shows that measurable trustworthiness characteristics (accuracy,
robustness, security, explainability, fairness) are broadly covered by established methods, that the
same evaluation artifact is often relevant to all three frameworks at once, and that process
obligations, human-oversight effectiveness, logging adequacy, and---most critically---the
behavior of agentic systems remain weakly served by current technical methods. We offered a
practitioner playbook for assembling a cross-framework evidence dossier and a research
agenda led by action-level evaluation for agents. Cited references and regulatory claims
were verified against official and primary sources where publicly accessible, and the
coding artifacts are released for external scrutiny. The mapping is a structured,
auditable argument graded by technical relevance---not a deterministic compliance tool; final
conformity determinations remain the responsibility of providers, deployers, and notified
bodies, and the mapping does not constitute legal advice. We release the
mapping as a living, machine-readable artifact in the hope that it becomes a durable
reference as both AI evaluation methods and AI regulation continue to mature.
